% 4. ผลการดำเนินงาน -- ย่อจากรายงานบทที่ 4 (../report/chapter/04-results.tex)
\section{ผลการดำเนินงาน}

\subsection{ลักษณะของกลุ่มตัวอย่างและคุณภาพข้อมูล}

ผู้ตอบ 107 คนเป็นหญิง 59 คน (ร้อยละ 55.1) และชาย 48 คน (ร้อยละ 44.9)
อายุเฉลี่ย 19.36 ปี (SD~= 0.82) ส่วนใหญ่เรียนชั้นปีที่ 2 (ร้อยละ 81.3)
แพลตฟอร์มที่ใช้มากที่สุดคือ TikTok (ร้อยละ 42.1) Instagram (ร้อยละ 27.1)
และ YouTube (ร้อยละ 25.2)
ประเภทเนื้อหาที่ดูมากที่สุดคือบันเทิง 58 คน รองลงมาคือชีวิตประจำวันของคนอื่น 13 คน
ส่วนข่าวสารบ้านเมืองและการเมืองมีเพียง 8 คน

ข้อมูลที่เก็บเองไม่มีลักษณะของข้อมูลสังเคราะห์
ชั่วโมงการใช้งานร้อยละ 90.7 ลงท้ายด้วย .0 หรือ .5 ซึ่งเป็นแบบที่คนจริงตอบ
ข้อมูลมีค่าว่าง 4 ช่องจาก 1,284 ช่อง
และสหสัมพันธ์ระหว่างตัวแปรต้นสูงสุดเพียง 0.45 (ชั่วโมงการนอนกับคุณภาพการนอน)
ขณะที่ชุดข้อมูลสังเคราะห์มีค่าสูงถึง 0.80 ถึง 0.95

ผู้ตอบใช้โซเชียลมีเดียเฉลี่ย 7.83 ชั่วโมงต่อวัน ไม่มีใครใช้น้อยกว่า 4 ชั่วโมง
(ตาราง~\ref{tab:variables}) และนอนเฉลี่ย 6.21 ชั่วโมงต่อคืน
ซึ่งน้อยกว่า 7 ชั่วโมงอย่างมีนัยสำคัญ (\stat{t}(105)~= −5.24, \stat{p}~< 0.001)
คะแนน DASS-21 ทั้งสามมิติเบ้ขวา ความวิตกกังวลเป็นมิติที่ผู้ตอบอยู่นอกระดับปกติมากที่สุด
คือ 73 คน (ร้อยละ 68.2) และร้อยละ 15.0 อยู่ในระดับรุนแรงมาก
ขณะที่ภาวะซึมเศร้าและความเครียดอยู่นอกระดับปกติร้อยละ 41.1 และ 32.7 ตามลำดับ
ตัวเลขนี้เป็นผลคัดกรองเบื้องต้น ไม่ใช่การวินิจฉัย
ค่าแอลฟาของครอนบาคของมิติภาวะซึมเศร้า ความวิตกกังวล และความเครียด
เท่ากับ 0.844, 0.737 และ 0.771 ตามลำดับ (ทั้งฉบับ 0.894) จึงผ่านเกณฑ์ 0.70 ทุกมิติ

\subsection{เวลาที่ใช้และการติดตามข่าวเชิงลบ}

ชั่วโมงใช้โซเชียลมีเดียต่อวันไม่สัมพันธ์กับมิติใดอย่างมีนัยสำคัญ
(ตาราง~\ref{tab:correlation} และภาพ~\ref{fig:correlation})
ค่า \stat{r} ทั้งสามค่าติดลบเล็กน้อย แต่ช่วงความเชื่อมั่นทุกช่วงคร่อมศูนย์
และ \stat{r}\textsuperscript{2} สูงสุดเพียง 0.022
หรือเวลาที่ใช้อธิบายความแปรปรวนของคะแนนได้ไม่ถึงร้อยละ 2.2
สหสัมพันธ์สเปียร์แมนให้ผลทางเดียวกัน ($\rho$ อยู่ระหว่าง −0.12 ถึง −0.03)
จำนวนวันที่ติดตามข่าวเชิงลบสัมพันธ์ทางบวกกับความวิตกกังวลมากกว่ามิติอื่น (\stat{r}~= 0.146)
ตามทิศทางที่คาดไว้ แต่ไม่มีนัยสำคัญ (\stat{p}~= 0.134)
แม้ควบคุมตัวแปรอื่นแล้วก็ยังไม่มีนัยสำคัญ (partial \stat{r}~= 0.166, \stat{p}~= 0.103)
ส่วนภาวะซึมเศร้าแทบไม่สัมพันธ์กับการติดตามข่าวเลย (\stat{r}~= −0.013)

\begin{table}[!t]
\centering\small
\caption{สหสัมพันธ์ระหว่างพฤติกรรมการใช้โซเชียลมีเดียกับคะแนน DASS-21}
\label{tab:correlation}
\begin{tabularx}{\columnwidth}{|L|c|c|}
\hline
\textbf{มิติ} & \textbf{\stat{r} (\stat{p})} & \textbf{partial \stat{r} (\stat{p})} \\
\hline
\multicolumn{3}{|l|}{\textit{ชั่วโมงใช้งานต่อวัน}} \\
\hline
ภาวะซึมเศร้า & −0.136 (0.161) & −0.101 (0.322) \\
ความวิตกกังวล & −0.148 (0.129) & −0.130 (0.200) \\
ความเครียด & −0.087 (0.376) & −0.064 (0.532) \\
\hline
\multicolumn{3}{|l|}{\textit{วันที่ติดตามข่าวเชิงลบ}} \\
\hline
ภาวะซึมเศร้า & −0.013 (0.892) & −0.001 (0.993) \\
ความวิตกกังวล & 0.146 (0.134) & 0.166 (0.103) \\
ความเครียด & 0.101 (0.303) & 0.106 (0.297) \\
\hline
\end{tabularx}

\vspace{2pt}
\raggedright หมายเหตุ: \stat{r} ใช้ $n = 107$ ส่วน partial \stat{r} ควบคุมชั่วโมงการนอน คุณภาพการนอน
ความกดดันจากการเรียน เกรดเฉลี่ย และความเพียงพอทางการเงิน ($n = 103$)
\end{table}

\begin{figure}[!t]
\centering
\includegraphics[width=\columnwidth]{fig-correlation.png}
\caption{ค่าสหสัมพันธ์ของชั่วโมงใช้งานและวันที่ติดตามข่าวเชิงลบกับคะแนน DASS-21
         เส้นแนวนอนคือช่วงความเชื่อมั่น 95\% ($n = 107$)}
\label{fig:correlation}
\end{figure}

\subsection{ประเภทเนื้อหาและการนอน}

คะแนนของผู้ที่ดูเนื้อหาต่างกลุ่มกันไม่ต่างกันในทุกมิติ
(Kruskal-Wallis \stat{p}~= 0.936, 0.883 และ 0.691
สำหรับภาวะซึมเศร้า ความวิตกกังวล และความเครียด)
กลุ่มข่าวมีคะแนนความวิตกกังวลเฉลี่ยสูงที่สุด (13.00) เทียบกับบันเทิง 11.59 อื่น~ๆ 11.00
และชีวิตคนอื่น 10.46 แต่ ANOVA ไม่พบความต่าง (\stat{F}(3,~103)~= 0.225, \stat{p}~= 0.879)
ทั้งนี้กลุ่มข่าวมีเพียง 8 คน จึงตรวจความต่างขนาดเล็กได้ยาก

ชั่วโมงใช้โซเชียลมีเดียไม่สัมพันธ์กับชั่วโมงการนอน (\stat{r}~= 0.049, \stat{p}~= 0.616)
และไม่สัมพันธ์กับคุณภาพการนอน (\stat{r}~= 0.114, \stat{p}~= 0.243)
แต่คุณภาพการนอนเองสัมพันธ์ทางลบกับภาวะซึมเศร้า (\stat{r}~= −0.276, \stat{p}~= 0.004)
และความเครียด (\stat{r}~= −0.320, \stat{p}~< 0.001)
ส่วนชั่วโมงการนอนและจำนวนวันที่ออกกำลังกายไม่สัมพันธ์กับมิติใดอย่างมีนัยสำคัญ

\subsection{เมื่อควบคุมปัจจัยอื่น}

ตัวแปรต้นทั้งแปดตัวมีค่า VIF ระหว่าง 1.05 ถึง 1.54
จึงไม่มีปัญหาตัวแปรต้นสัมพันธ์กันเอง
ตาราง~\ref{tab:regression} แสดงว่าตัวแปรโซเชียลมีเดียทั้งสองตัวเพิ่ม \Rsq{} ได้น้อย
และไม่มีนัยสำคัญในมิติใด
มิติความวิตกกังวลเพิ่มมากที่สุด ($\Delta$\Rsq~= 0.045, \stat{p}~= 0.100)
และหลังปรับ Bonferroni ค่า \stat{p} เป็น 0.301

\begin{table}[!t]
\centering\small
\caption{สัมประสิทธิ์ \stat{b} ของการถดถอยพหุคูณแบบลำดับขั้น ($n = 103$)}
\label{tab:regression}
\begin{tabularx}{\columnwidth}{|L|r|r|r|}
\hline
\textbf{ตัวแปรต้น} & \textbf{ซึมเศร้า} & \textbf{วิตกกังวล} & \textbf{เครียด} \\
\hline
ชั่วโมงใช้งานต่อวัน & −0.268 & −0.391 & −0.187 \\
วันที่ติดตามข่าวเชิงลบ & 0.023 & 0.872 & 0.534 \\
ชั่วโมงการนอน & 1.150 & 0.069 & 0.195 \\
คุณภาพการนอน & −0.921* & −0.162 & −0.878* \\
วันที่ออกกำลังกาย & 0.176 & −0.066 & −0.185 \\
ความกดดันจากการเรียน & 0.979** & 0.329 & 0.800* \\
เกรดเฉลี่ยสะสม & −0.706 & −0.958 & 0.115 \\
ความเพียงพอทางการเงิน & −1.125** & −0.663* & −0.626 \\
\hline
\Rsq{} ขั้นที่ 1 & 0.263 & 0.058 & 0.191 \\
\Rsq{} ขั้นที่ 2 & 0.271 & 0.103 & 0.205 \\
$\Delta$\Rsq & 0.008 & 0.045 & 0.013 \\
\stat{p} ของ $\Delta$\Rsq & 0.615 & 0.100 & 0.454 \\
\hline
\end{tabularx}

\vspace{2pt}
\raggedright หมายเหตุ: * \stat{p}~< 0.05, ** \stat{p}~< 0.01 ก่อนปรับ Bonferroni
ค่าความคลาดเคลื่อนมาตรฐานแบบ HC3
\end{table}

ตัวแปรที่มีน้ำหนักในสมการเป็นตัวแปรควบคุม ไม่ใช่ตัวแปรโซเชียลมีเดีย
ในสมการภาวะซึมเศร้า ความเพียงพอทางการเงินมีน้ำหนักมากที่สุด ($\beta = -0.293$)
ตามด้วยความกดดันจากการเรียน ($\beta = 0.269$) และคุณภาพการนอน ($\beta = -0.242$)
เมื่อปรับ Bonferroni แล้ว เหลือเพียงความเพียงพอทางการเงิน (\stat{p}~= 0.002)
และความกดดันจากการเรียน (\stat{p}~= 0.026) ในสมการภาวะซึมเศร้าที่ยังมีนัยสำคัญ
ในสมการความเครียดซึ่งควบคุมเกรดเฉลี่ยและสภาพการเงินแล้ว
สัมประสิทธิ์ของชั่วโมงใช้งานคือ \stat{b}~= −0.187 คะแนนต่อชั่วโมง
(95\% CI −0.71 ถึง 0.34, \stat{p}~= 0.485)
ผู้ที่ใช้โซเชียลมีเดียมากกว่าจึงไม่ได้เครียดต่างจากผู้ที่ใช้น้อยกว่า
การตรวจข้อตกลงเบื้องต้นพบความแปรปรวนไม่คงที่ในสมการภาวะซึมเศร้า
(Breusch-Pagan \stat{p}~= 0.013) ซึ่งเป็นเหตุผลที่ใช้ HC3
และค่าคลาดเคลื่อนของสมการความวิตกกังวลไม่แจกแจงปกติ (Shapiro-Wilk \stat{p}~= 0.001)
จากผู้ตอบไม่กี่คนที่มีคะแนนสูงมาก

\subsection{สภาพการเงิน}

เมื่อแบ่งผู้ตอบตามคะแนนความเพียงพอทางการเงินเป็น 3 กลุ่ม คือ ไม่พอใช้ (1--4) พอใช้ (5--7)
และสบาย (8--10) คะแนนภาวะซึมเศร้าต่างกันมากที่สุด (ภาพ~\ref{fig:finance})
กลุ่มไม่พอใช้มีคะแนนเฉลี่ย 16.6 ขณะที่กลุ่มสบายมี 6.6
(\stat{F}(2,~104)~= 10.91, \stat{p}~< 0.001, \etasq{}~= 0.173, Cohen's \stat{d}~= 1.59)
ความเครียดต่างกันเช่นกัน (\stat{F}(2,~104)~= 5.02, \stat{p}~= 0.008, \etasq{}~= 0.088)
ส่วนความวิตกกังวลไม่ต่างกัน (\stat{p}~= 0.153) และการทดสอบ Kruskal-Wallis ให้ข้อสรุปเดียวกัน
เมื่อใช้คะแนน 1--10 โดยตรง ความเพียงพอทางการเงินสัมพันธ์กับภาวะซึมเศร้า (\stat{r}~= −0.374)
แรงกว่าชั่วโมงใช้โซเชียลมีเดีย (\stat{r}~= −0.136) เกือบสามเท่า
ส่วนเกรดเฉลี่ยสะสมไม่สัมพันธ์กับมิติใด ($|r| \le 0.063$)

\begin{figure}[!t]
\centering
\includegraphics[width=\columnwidth]{fig-finance.png}
\caption{คะแนนเฉลี่ย DASS-21 แยกตามความเพียงพอทางการเงิน ไม่พอใช้ (1--4) พอใช้ (5--7)
         และสบาย (8--10) เส้นบนแท่งคือค่าคลาดเคลื่อนมาตรฐาน}
\label{fig:finance}
\end{figure}

\subsection{ผลเมื่อรวมผู้ที่ไม่ผ่านข้อดักความใส่ใจ}

เมื่อวิเคราะห์ชุดเดิมกับผู้ตอบทั้ง 121 คน ข้อสรุปหลักยังเหมือนเดิม แต่มีผลสองข้อที่เปลี่ยน
คือ สหสัมพันธ์บางส่วนระหว่างการติดตามข่าวกับความวิตกกังวลกลายเป็นมีนัยสำคัญ
(partial \stat{r}~= 0.189, \stat{p}~= 0.047)
และคะแนนภาวะซึมเศร้าของชายและหญิงต่างกันอย่างมีนัยสำคัญ (\stat{p}~= 0.049)
ผลทั้งสองข้อขึ้นกับผู้ตอบที่ไม่ได้อ่านคำถามจริง จึงไม่นำมาสรุป
